\documentclass[10pt,letterpaper]{article}
\usepackage[margin=0.78in]{geometry}
\usepackage{booktabs}
\usepackage{xcolor}
\usepackage[hidelinks]{hyperref}
\usepackage{array}
\pagestyle{empty}
\setlength{\parindent}{0pt}
\setlength{\parskip}{5pt}
\definecolor{riceblue}{HTML}{124E78}
\begin{document}
{\Large\bfseries Sound Calibration: Bench Report}\\[-2pt]
{\small Rice Auditory Benchmarking \quad | \quad 28 October 2026}

\textbf{Purpose.} Make the measured sound components of 14 short tone pips more consistent across frequency. This is a bench test, not a validated animal stimulus.

\textbf{Setup and method.} An Audient iD14 played reconstructed 10 ms pips with 2 ms ramps at 96 kHz; an Earthworks M50 recorded the sound. Howard reported the microphone about 10 cm from the speaker, then moved it sideways twice by about 5 mm to improve high-frequency pickup. The final correction applies to that position and the same gain settings. A REED R8090 calibrator, reported set to 114 dB and fitted using an adapter, supplied a near-1 kHz reference.

The software reduces the digital level of stronger frequencies to bring their measured microphone components closer together; no frequency was boosted above the $-36$ dBFS base. The numbers below are \emph{not} the digital playback settings.

\textbf{Results.} Two seven-second tests each presented 700 consecutive pips. The spread across the 14 frequency-group medians was \textbf{1.13 dB} and \textbf{0.77 dB}. These compare the recorded frequency component in the middle 6 ms, not every individual pip.

\begin{center}
\small
\begin{tabular}{@{}rr@{\hspace{1.8em}}rr@{}}
\toprule
Frequency & Estimate & Frequency & Estimate\\
(kHz) & (dB) & (kHz) & (dB)\\
\midrule
2    & 41.3 & 18.2 & 41.6\\
4    & 42.1 & 19.5 & 41.8\\
6    & 41.6 & 20.9 & 41.6\\
8    & 41.6 & 22.4 & 41.7\\
9    & 41.7 & 24   & 41.7\\
10.5 & 41.9 & 25.8 & 41.7\\
17   & 41.6 & 27   & 41.9\\
\bottomrule
\end{tabular}
\end{center}
{\footnotesize\textit{Table:} Repeat-run frequency-group median estimates for the central component, expressed as equivalent-sine levels relative to 20 $\mu$Pa. Values assume the nominal 114 dB reference, unchanged linear gain and flat microphone response. Decimal places aid comparison; accuracy is not established.}

\textbf{Limits and next steps.} These estimates are \textbf{not validated whole-pip dB SPL or standardized peak-equivalent SPL}. Background sound was substantial in full 10 ms windows; the microphone's individual electronic response, distortion and measurement uncertainty have not been verified. Equal measured pressure also does not establish equal perceived loudness. Recheck the reference, define the pip-level measurement, and repeat at the intended sheep-ear position before experimental use.

{\footnotesize Evidence: \href{https://howardwhsrun.github.io/audiotory_benchmark/web/calibration/}{Sound calibration page} and \href{https://github.com/HowardWHSrun/audiotory_benchmark/blob/main/docs/verification/2026-10-28_Conditional_Pip_SPL_Estimates.md}{measurement record}.}
\end{document}